\documentclass[tikz,border=8pt]{standalone}
\usepackage{ctex}
\usepackage{amsmath}
\usepackage{pgfplots}
\pgfplotsset{compat=1.18}
\usetikzlibrary{calc,positioning,arrows.meta}

%% 图 opt-p2-06-1：牛顿法 vs 梯度下降轨迹对比
%% 相同椭圆等高线：GD 之字形多步 / 牛顿一步直达
%% 说明 Hessian 信息加速

\begin{document}
\begin{tikzpicture}[scale=1.2]

  % 坐标轴
  \draw[->, gray!55, thin] (-3.2, 0) -- (3.2, 0)
    node[right, font=\scriptsize] {$x_1$};
  \draw[->, gray!55, thin] (0, -2.5) -- (0, 2.8)
    node[above, font=\scriptsize] {$x_2$};

  % 椭圆等高线（f = x1^2/4 + x2^2，kappa=4，横轴长是纵轴2倍）
  \foreach \r in {0.5, 0.9, 1.3, 1.7, 2.1}{
    \draw[gray!40, thin] (0,0) ellipse ({2*\r} and {\r});
  }
  \node[font=\scriptsize, gray!60] at (1.2, 0.45) {$c_1$};
  \node[font=\scriptsize, gray!60] at (2.05, 0.8) {$c_2$};

  % 极小值点
  \fill[black] (0,0) circle (2.5pt)
    node[below right, font=\small] {$\mathbf{x}^*$};

  % ---- 梯度下降轨迹（红色，之字形）----
  % 从 x0=(2.2, 1.0) 出发，步长 alpha=0.4
  % 梯度 = (x1/2, 2*x2)，模拟之字形
  \coordinate (G0) at (2.2, 1.0);
  \coordinate (G1) at (1.75, -0.78);
  \coordinate (G2) at (1.22, 0.58);
  \coordinate (G3) at (0.85, -0.40);
  \coordinate (G4) at (0.55, 0.26);
  \coordinate (G5) at (0.28, -0.13);

  \draw[->, red!75!black, thick, >=Stealth] (G0)--(G1);
  \draw[->, red!75!black, thick, >=Stealth] (G1)--(G2);
  \draw[->, red!75!black, thick, >=Stealth] (G2)--(G3);
  \draw[->, red!75!black, thick, >=Stealth] (G3)--(G4);
  \draw[->, red!75!black, thick, >=Stealth] (G4)--(G5);
  \draw[->, red!75!black, thick, dashed, >=Stealth] (G5) -- (0.10, -0.05);

  \fill[red!75!black] (G0) circle (2.5pt)
    node[above right, font=\small] {$\mathbf{x}_0$};
  \foreach \pt in {G1, G2, G3, G4, G5}{
    \fill[red!75!black] (\pt) circle (2pt);
  }
  \node[font=\scriptsize, red!75!black] at (G1) [right=2pt] {$\mathbf{x}_1$};

  % GD 图例标注
  \node[font=\small, red!75!black, fill=white, inner sep=2pt]
    at (2.5, 1.8)
    {梯度下降};
  \draw[red!75!black, thick] (1.8, 1.75) -- (2.1, 1.75);

  % ---- 牛顿法轨迹（蓝色，一步直达）----
  % 从同一初始点 x0=(2.2, 1.0) 出发
  % H^{-1} = diag(2, 1/2)，grad f = (1.1, 2.0)
  % d = -H^{-1}*g = -(2.2, 1.0) = (-2.2, -1.0)（二次函数精确一步）
  \coordinate (N0) at (2.2, 1.0);

  \draw[->, blue!80!black, very thick, >=Stealth] (N0) -- (-0.02, -0.01);

  \fill[blue!80!black] (N0) circle (2.5pt);
  \node[font=\scriptsize, blue!80!black] at (1.4, 0.7) {牛顿一步};
  % 箭头指向路径中点
  \draw[->, blue!70!black, thin] (1.4, 0.65) -- (1.15, 0.53);

  % 牛顿图例标注
  \node[font=\small, blue!80!black, fill=white, inner sep=2pt]
    at (2.5, 2.2)
    {牛顿法（$\alpha=1$）};
  \draw[blue!80!black, very thick] (1.8, 2.15) -- (2.1, 2.15);

  % ---- Hessian 信息说明框 ----
  \node[draw=gray!50, thin, fill=white, rounded corners=3pt,
        font=\scriptsize, inner sep=5pt, align=left]
    at (-1.8, -1.8)
    {梯度下降：方向 $= -\nabla f$\\
     只用一阶信息，受条件数影响\\[3pt]
     牛顿法：方向 $= -\mathbf{H}^{-1}\nabla f$\\
     Hessian 补偿尺度差异，仿射不变};

  % 整体标题
  \node[font=\large\bfseries] at (0, 3.2)
    {牛顿法 vs 梯度下降：相同等高线，不同轨迹};

\end{tikzpicture}
\end{document}
